% !TEX root = ../Thesis.tex

\chapter{Definizioni e terminologia}
\label{ch:definizioni}
\setlist[itemize, 1]{label= , leftmargin=*, itemsep=0pt}

Nel seguito sono presentati un insieme di termini e definizioni fondamentali per la comprensione dei concetti chiave trattati nella tesi. I termini sono organizzati per aree tematiche omogenee, procedendo dai fondamenti della comunicazione sicura fino agli aspetti più specialistici della generazione steganografica. Questo lessico tecnico costituisce il vocabolario essenziale per navigare efficacemente all'interno del discorso sulla steganografia generativa, sui modelli linguistici, sulle architetture di reti neurali e sui generatori di immagini.

\section{Concetti fondazionali della comunicazione}
\begin{itemize}
	\item \textbf{Plaintext}: informazione nella sua forma originaria e integralmente leggibile, antecedente a qualsiasi trasformazione crittografica o steganografica. Nel contesto della presente trattazione, il plaintext costituisce il dato in chiaro che si intende proteggere.
	\item \textbf{Ciphertext}: risultato dell'applicazione di un algoritmo di cifratura a un plaintext. Tale rappresentazione, incomprensibile in assenza della corrispondente chiave di decifrazione, costituisce il prodotto finale del processo crittografico.
	\item \textbf{Coverobject}: oggetto multimediale originale, nella sua forma inalterata, che funge da contenitore all'interno del quale si intende occultare un messaggio segreto.
	\item \textbf{Stegoobject}: oggetto risultante dall'incorporazione del messaggio segreto all'interno del coverobject. La proprietà fondamentale che tale artefatto deve soddisfare è l'indistinguibilità percettiva e statistica dal coverobject originale. Nel caso di \emph{Meteor} o delle VQ-GAN, il termine si riferisce all'oggetto \textbf{generato} che contiene l'informazione segreta codificata, che non è in senso stretto un "oggetto modificato" ma un "oggetto creato" ad hoc. In questo contesto, il termine "stegoobject" è usato in modo estensivo per indicare qualsiasi output generato che contenga un messaggio segreto, indipendentemente dal fatto che sia stato ottenuto tramite embedding o generazione ex novo.
	\item \textbf{Payload}: quantità di informazione segreta, espressa in bit, che può essere occultata all'interno di un coverobject senza comprometterne l'integrità statistica e percettiva.
\end{itemize}

\section{Discipline della comunicazione sicura}
\begin{itemize}
	\item \textbf{Crittografia}: disciplina formale che si occupa di rendere non intelligibile un messaggio ad attori non autorizzati mediante l'applicazione di metodi di cifratura. La crittografia costituisce la tecnica preminente per garantire la confidenzialità dei dati nei sistemi informatici contemporanei. Tale disciplina si fonda su principi matematici e algoritmi atti a trasformare l'informazione in una forma cifrata, decifrabile esclusivamente da chi possiede la chiave corretta. Pur costituendo il fondamento dei sistemi comunicativi moderni, la crittografia è potenzialmente vulnerabile ad attacchi di forza bruta, a exploit basati su vulnerabilità specifiche degli algoritmi, a compromissioni delle chiavi e, in prospettiva, ad algoritmi di decriptazione quantistica. Quando implementata correttamente, la crittografia fornisce un livello di sicurezza elevato, rendendo computazionalmente proibitivo l'accesso non autorizzato ai dati \cite{neskey-2025}.
	\item \textbf{Cifratura simmetrica}: paradigma crittografico in cui il medesimo segreto condiviso viene impiegato tanto per la cifratura quanto per la decifrazione del messaggio. Tale metodologia, generalmente caratterizzata da una maggiore efficienza computazionale rispetto alla cifratura asimmetrica, richiede che la chiave sia trasmessa in modo sicuro tra le parti coinvolte.
	\item \textbf{Steganografia}: disciplina che si occupa di occultare l'esistenza stessa di un messaggio all'interno di un oggetto multimediale, rendendolo impercettibile a osservatori non autorizzati. La steganografia si configura come tecnica complementare alla crittografia: mentre quest'ultima protegge il contenuto informativo, la steganografia tutela la riservatezza dell'atto comunicativo. Storicamente impiegata per finalità di spionaggio e comunicazione clandestina, la steganografia ha trovato nuove applicazioni nell'era digitale, quali la protezione dei diritti d'autore, la comunicazione sicura e la salvaguardia della privacy in rete \cite{turch-20}. Le modalità di implementazione includono l'inserimento di dati occulti in immagini, tracce audio o flussi video, nonché l'impiego di tecniche di codifica testuale. La steganografia è tuttavia vulnerabile a tecniche di rilevamento avanzate (steganalisi) che analizzano le proprietà statistiche degli oggetti multimediali per identificare anomalie indicative della presenza di dati nascosti.
	\item \textbf{Steganalisi}: disciplina deputata al rilevamento di messaggi occulti all'interno di oggetti multimediali, avvalendosi di metodologie di analisi statistica e di apprendimento automatico. La steganalisi costituisce il corrispettivo avversario della steganografia, concentrandosi sull'individuazione dell'esistenza stessa del messaggio celato piuttosto che sulla sua protezione. Con l'avvento delle tecnologie digitali, la steganalisi ha raggiunto livelli di sofisticazione notevoli, impiegando algoritmi avanzati per l'esame delle caratteristiche statistiche degli oggetti multimediali \cite{apau-24}.
	\item \textbf{Steganografia senza embedding (SwE, Steganography without Embedding)}: paradigma steganografico in cui il messaggio segreto non viene incorporato (\emph{embedded}) all'interno di un coverobject preesistente, bensì il coverobject stesso viene \textbf{generato ex novo} a partire dal messaggio segreto. In questo approccio, un modello generativo (quale un modello linguistico o un generatore di immagini) produce un artefatto multimediale -- testo, immagine o audio -- la cui struttura e contenuto sono determinati dai bit del messaggio da occultare. Poiché non esiste un coverobject originale di riferimento, la steganalisi classica basata sul confronto tra coverobject e stegoobject risulta inefficace. L'indistinguibilità è garantita dalla qualità del modello generativo, che deve produrre output statisticamente e percettivamente equivalenti a campioni naturali. Paradigmi come \emph{Meteor} per rappresentano istanze di questo approccio.
\end{itemize}

\section{Fondamenti di modellazione linguistica}
\begin{itemize}
	\item \textbf{Modelli linguistici}: classe di modelli di intelligenza artificiale progettati per comprendere e generare testo in linguaggio naturale. Tali modelli vengono addestrati su ingenti quantità di dati testuali e impiegano tecniche di apprendimento automatico per apprendere le regolarità grammaticali, il significato delle parole e le relazioni sintattico-semantiche tra le frasi. I modelli linguistici trovano applicazione in sistemi quali chatbot, piattaforme di traduzione automatica, generatori di testo e analizzatori del sentiment. Essi possono presentare vulnerabilità riconducibili a bias nei dati di addestramento e a limitazioni nella comprensione del contesto.
	\item \textbf{Modelli di linguaggio di grandi dimensioni (LLM)}: sottocategoria di modelli linguistici caratterizzata da un numero di parametri elevato, che consente di apprendere rappresentazioni più complesse del linguaggio naturale e di generare testo con maggiore coerenza e naturalezza. Gli LLM vengono addestrati su corpora testuali di dimensioni massive e si avvalgono di architetture avanzate, quali i Transformer, per catturare le dipendenze a lungo termine tra le unità linguistiche.
	\item \textbf{Modelli generativi}: classe di modelli di intelligenza artificiale atti a generare nuovi dati che presentano caratteristiche analoghe a quelli impiegati in fase di addestramento. Attraverso tecniche di apprendimento automatico, tali modelli apprendono la distribuzione sottostante ai dati di training e producono nuovi campioni che seguono la medesima distribuzione. I modelli generativi trovano impiego in numerosi ambiti, inclusi la generazione testuale, la sintesi di immagini, la composizione musicale e la sintesi vocale.
\end{itemize}

\section{Componenti dei modelli linguistici}
\begin{itemize}
	\item \textbf{Token}: unità testuale elementare impiegata nei modelli linguistici, corrispondente a una parola, un carattere o una sottoporzione di parola. I token costituiscono l'unità fondamentale su cui operano i modelli linguistici, rappresentando le entità minime che il modello deve elaborare e generare. La tokenizzazione -- il processo di suddivisione del testo in token -- può essere effettuata secondo diverse modalità, quali la segmentazione per parole, per caratteri o per sottoparole (ad esempio mediante algoritmi \emph{subword} come Byte Pair Encoding).
	\item \textbf{Vocabolario}: insieme dei token che un modello linguistico è in grado di processare e generare. Il vocabolario viene definito durante la fase di addestramento e la sua dimensione può variare in funzione del modello e dell'applicazione. Un vocabolario più esteso consente al modello di coprire una gamma più ampia di espressioni linguistiche, ma comporta un aumento della complessità computazionale.
	\item \textbf{Logits}: valori numerici prodotti dall'ultimo strato di un modello generativo prima dell'applicazione della funzione di attivazione (tipicamente la softmax) per ottenere una distribuzione di probabilità sui token successivi. I logits rappresentano le preferenze del modello per ciascun token del vocabolario, dove valori più elevati indicano una maggiore propensione alla selezione.
	\item \textbf{Softmax}: funzione di attivazione impiegata nei modelli linguistici per trasformare il vettore dei logits in una distribuzione di probabilità sull'insieme dei token del vocabolario. Formalmente, $\mathit{softmax}(x_i) = e^{x_i} / \sum_{j} e^{x_j}$, dove $x_i$ rappresenta il logit corrispondente al token $i$-esimo e la somma al denominatore è estesa all'intero vocabolario.
\end{itemize}

\section{Meccanismi di generazione e controllo}
\begin{itemize}
	\item \textbf{Prompt}: sequenza testuale iniziale fornita in input a un modello linguistico per orientare e condizionare la generazione del testo successivo, definendo il contesto semantico e tematico dell'output prodotto.
	\item \textbf{Contesto}: unione dei token del prompt e dei token generati fino a un dato punto temporale, che costituisce l'input su cui il modello linguistico basa la generazione del token successivo. Il contesto è dinamico e si evolve iterativamente con l'aggiunta di nuovi token, influenzando le distribuzioni di probabilità generate dal modello a ogni passo.
	\item \textbf{Finestra di contesto}: limite massimo di token che un modello può considerare come input per la generazione del token successivo. La dimensione della finestra di contesto influisce sulla capacità del modello di mantenere coerenza e riferimenti a lungo termine all'interno del testo generato. 
	\item \textbf{Entropia}: nel contesto attuale, misura della variabilità o incertezza associata alla distribuzione di probabilità dei token successivi generata da un modello linguistico. Un'alta entropia indica una distribuzione pi\`u uniforme, con molte opzioni possibili per il token successivo, mentre una bassa entropia suggerisce una distribuzione concentrata su pochi token altamente probabili. L'entropia è un fattore cruciale nella steganografia generativa, poiché determina la capacità del sistema di codificare informazione senza compromettere la naturalezza del testo prodotto.
	\item \textbf{Campionamento}: processo di selezione di un token successivo a partire da una distribuzione di probabilità generata da un modello. Le modalità di campionamento includono il campionamento casuale (\emph{random sampling}), la selezione deterministica del token più probabile (\emph{greedy decoding}), il campionamento casuale solo sui $k$ token più probabili (\emph{top-k sampling}), il campionamento casuale solo sui token con probabilità sopra una certa soglia (\emph{top-p sampling} o \emph{nucleus sampling}) e il campionamento vincolato, come nel caso di \emph{Meteor}, in cui la scelta del token è determinata dai bit del messaggio segreto.
\end{itemize}

\section{Architetture di reti neurali}
\begin{itemize}
	\item \textbf{Transformer}: architettura di rete neurale basata sul meccanismo di attenzione (\emph{self-attention}), ampiamente adottata nei modelli linguistici di grandi dimensioni per la capacità di catturare relazioni a lungo termine tra gli elementi di una sequenza, superando le limitazioni delle architetture ricorrenti nella modellazione di dipendenze distanti.
	\item \textbf{Convolutional layer}: strato che applica operazioni di convoluzione per estrarre caratteristiche locali dai dati in ingresso, comunemente impiegato nelle reti neurali convoluzionali (CNN) per l'elaborazione di segnali strutturati su griglia, quali immagini e volumi multidimensionali.
	\item \textbf{Recurrent layer}: strato che introduce connessioni ricorrenti per mantenere e propagare informazioni sullo stato precedente della sequenza, tipicamente utilizzato nelle reti neurali ricorrenti (RNN) per l'elaborazione di dati sequenziali, quali testo, serie temporali e segnali audio.
	\item \textbf{Fully connected layer}: strato in cui ogni neurone è connesso a tutti i neuroni dello strato precedente, impiegato principalmente per la classificazione o la regressione dopo l'estrazione di caratteristiche da parte di strati convoluzionali o ricorrenti.
\end{itemize}

\section{Concetti fondamentali della rappresentazione visiva}
\begin{itemize}
	\item \textbf{Patch}: porzione di un'immagine, tipicamente di dimensioni fisse, che viene elaborata come unità discreta in modelli di generazione visuale. I patch consentono di suddividere un'immagine in segmenti più gestibili, facilitando l'estrazione di caratteristiche locali e la generazione di contenuti visivi coerenti.
	\item \textbf{Spazio latente}: rappresentazione astratta e compressa dei dati in ingresso, appresa da un modello generativo durante la fase di addestramento. Lo spazio latente consente al modello di catturare le caratteristiche essenziali dei dati, facilitando la generazione di nuovi campioni che condividono le stesse proprietà statistiche e semantiche degli esempi di training. Per spazio latente si intende un dominio di rappresentazione in cui i dati originali vengono mappati sulla base delle feature, consentendo al modello di manipolare e generare nuovi campioni attraverso operazioni su queste rappresentazioni astratte. Si utilizza la distanza euclidea o altre metriche di similarità per misurare la vicinanza tra punti nello spazio latente, che riflette la \emph{somiglianza semantica o visiva} tra i dati.
	\item \textbf{Semantica di un'immagine}: insieme di significati, concetti e relazioni rappresentati visivamente in un'immagine. La semantica di un'immagine comprende elementi quali oggetti, scene, azioni e contesti che possono essere interpretati da un osservatore umano o da un modello di intelligenza artificiale. La comprensione della semantica di un'immagine è fondamentale per applicazioni come il riconoscimento delle immagini, la generazione di descrizioni testuali e la manipolazione visiva.
	\item \textbf{Sintassi di un'immagine}: insieme di regole e strutture che governano l'organizzazione visiva degli elementi all'interno di un'immagine. La sintassi di un'immagine riguarda aspetti quali la composizione, la disposizione spaziale, la prospettiva e le relazioni formali tra gli oggetti rappresentati. La sintassi contribuisce alla percezione estetica e alla comprensione visiva, influenzando il modo in cui un'immagine viene interpretata e apprezzata.
\end{itemize}

\section{Architetture per la generazione visuale}
\begin{itemize}
	\item \textbf{Generative Adversarial Network (GAN)}: architettura di rete neurale composta da due componenti distinti -- un generatore e un discriminatore -- che competono tra loro durante la fase di addestramento. Il generatore tenta di produrre dati sintetici indistinguibili da quelli reali, mentre il discriminatore si adopera per discriminare tra campioni autentici e campioni generati. Questo processo di apprendimento avversario conduce a un miglioramento progressivo della qualità dei dati sintetici prodotti.
	\item \textbf{Vector Quantized-Generative Adversarial Network (VQ-GAN)}: architettura di rete neurale che integra i principi dei GAN con tecniche di quantizzazione vettoriale, impiegando un codebook discreto per rappresentare le caratteristiche latenti dei dati in ingresso. La VQ-GAN è progettata per migliorare la qualità e la stabilità del processo generativo, riducendo al contempo la complessità computazionale mediante la discretizzazione dello spazio latente \cite{esser-2020}.
	\item \textbf{Codebook}: insieme strutturato di vettori latenti appresi durante l'addestramento di un modello basato su quantizzazione vettoriale, come la VQ-GAN. Ciascun vettore del codebook costituisce un elemento atomico del dizionario di rappresentazioni discrete mediante il quale i dati originali vengono codificati e successivamente ricostruiti.
	\item \textbf{Quantizzazione}: processo di discretizzazione di uno spazio di rappresentazione continuo (o discreto ma di dimensione eccessivamente elevata) in un insieme finito di valori atomici detti \emph{token}. Tale procedura viene impiegata in architetture come la VQ-GAN per ridurre la complessità del modello e migliorare la stabilità del processo di addestramento, trasformando lo spazio latente continuo in un codebook finito di vettori discreti.
\end{itemize}

\section{Misure e garanzie nei sistemi steganografici}
\begin{itemize}
	\item \textbf{Bitrate variabile}: strategia di codifica in cui la quantità di informazione incorporata in ciascun token può variare dinamicamente in funzione delle caratteristiche del contesto e della distribuzione di probabilità dei token successivi. In contesti a bassa entropia, il sistema può astenersi dalla codifica di informazione, mentre in scenari ad alta entropia può codificare un numero maggiore di bit per token, massimizzando il throughput informativo senza compromettere la naturalezza statistica del testo generato.
	\item \textbf{Throughput informativo}: quantità di informazione segreta, misurata in bit, che può essere codificata e trasmessa attraverso un canale di comunicazione, quale un testo generato da un modello linguistico. Il throughput informativo è determinato da molteplici fattori, tra cui la capacità di codifica del modello sottostante, la distribuzione di probabilità dei token e la strategia di codifica adottata.
	\item \textbf{Anomalia semantica}: incoerenza o incongruenza nel significato di un testo generato, riconducibile a selezioni di token non naturali o a errori nel processo generativo. Le anomalie semantiche costituiscono indicatori potenzialmente rivelatori per strumenti di steganalisi, in quanto segnalano una deviazione dal comportamento generativo naturale.
	\item \textbf{Soglia di entropia minima}: valore critico dell'entropia al di sotto del quale un sistema di steganografia generativa, quale \emph{Meteor}, si astiene dalla codifica di informazione segreta. Tale soglia garantisce che l'inserimento di informazioni occulte avvenga esclusivamente in contesti in cui la distribuzione di probabilità dei token successivi presenta una varianza sufficiente, prevenendo la generazione di anomalie semantiche.
	\item \textbf{Soglia di probabilità}: valore prefissato impiegato per troncare la distribuzione di probabilità dei token successivi prodotta da un modello linguistico, eliminando i token la cui probabilità risulta inferiore a tale soglia. Questo meccanismo, noto come troncamento della coda (\emph{tail truncation}), circoscrive le scelte ai token maggiormente probabili, migliorando la coerenza e la naturalezza del testo generato.
	\item \textbf{Troncamento della coda}: tecnica impiegata nei modelli linguistici per limitare l'insieme dei token candidati alla selezione a quelli la cui probabilità supera una soglia prefissata. Tale procedura migliora la coerenza e la naturalezza del testo generato, eliminando i token meno probabili che potrebbero introdurre anomalie semantiche o statistiche. Questa tecnica riveste particolare rilevanza in sistemi di steganografia generativa, come \emph{Meteor}, in quanto contribuisce a garantire che il testo prodotto rimanga statisticamente indistinguibile da un testo naturale.
	\item \textbf{Indistinguibilità statistica}: proprietà di un testo generato da un modello linguistico, quale \emph{Meteor}, per cui le sue caratteristiche statistiche risultano indistinguibili da quelle di un testo naturale non steganografico. Tale proprietà è fondamentale per impedire a strumenti di steganalisi di rilevare la presenza di informazioni occulte basandosi su anomalie statistiche.
	\item \textbf{Indistinguibilità percettiva}: propriet\`a di un testo generato da un modello linguistico, quale \emph{Meteor}, per cui il suo significato e la sua coerenza risultano indistinguibili da quelli di un testo naturale non steganografico. Tale propriet\`a contribuisce a garantire che il testo prodotto non susciti sospetti basati su anomalie semantiche.
	\item \textbf{Non-determinismo}: caratteristica di un sistema di generazione per cui l'output prodotto non è deterministico ma dipende da fattori probabilistici, tra cui la distribuzione di probabilit\`a dei token successivi e la strategia di codifica steganografica adottata. Il non-determinismo contribuisce a garantire l'indistinguibilità statistica e percettiva del messaggio generato. Nel sistema implementato in questa tesi, il non-determinismo ha tuttavia rappresentato una criticità: l'architettura VQ-GAN impiega un codificatore (\emph{Encoder}) e un decodificatore (\emph{Decoder}) che, essendo funzioni statistiche apprese e non inversi esatti, introducono una componente d'errore nel processo di ricostruzione. Questa mancata biettività ha prodotto inconsistenze tra la codifica e la decodifica di alcune patch dell'immagine.
\end{itemize}

